\documentclass{article}
\usepackage{pgfplots}
\usepgfplotslibrary{external}
\usepackage[dvipsnames]{xcolor}
\tikzexternalize
\usepackage{fullpage}
\usepackage{ulem}
\usepackage{subcaption}
\usepackage{amssymb, amsmath, changepage,pgfplots, darkmode,tabularx,siunitx}
\usepackage{algorithm}
\enabledarkmode
\usepackage{makecell}
\usepackage{amsfonts}
\pgfplotsset{compat = newest}


\usepackage{tikz}
\usetikzlibrary{shadings}
\usetikzlibrary{calc,arrows.meta}
\usepackage[most]{tcolorbox}
\definecolor{pag}{HTML}{293133}
	\usepgfplotslibrary{colormaps,patchplots}
\color{white}

\tikzset{basic/.style={draw,text width=1em,text badly centered}}
\tikzset{input/.style={basic,circle}}
\tikzset{weights/.style={basic,rectangle}}
\tikzset{functions/.style={basic,circle}}

\title{Machine Learning Notes \\ Recurrent Neural Networks}
\author{Vitor Negromonte Cabral de Oliveira}

\date{June 2025}

\begin{document}

\maketitle

\newpage
\tableofcontents
\newpage
\section{Introdução à Redes Neurais Recorrentes}

Redes neurais tradicionais (MLPs), funcionam, de modo geral, como uma rede de \textit{um para um}, onde, por exemplo, envia-se uma imagem para o modelo e recebe-se um vetor de probabilidades para aquela imagem. Não obstante, é notório que essa situação é limitante quanto ao funcionamento de sistemas dessa natureza para demais cenários, onde, por exemplo, precisamos processar sequências de determinada natureza.

Neste cenário, temos as redes neurais recorrentes. Uma rede recorrente possui um estado interno e recebe um vetor $x_t$, ou seja, um vetor $x$ em um momento qualquer $_t$ e modifica o seu estado interno com uma função da entrada $x$ ao longo do tempo.

Matematicamente falando, uma rede recorrente modela o estado atual de um sistema como uma função $f_{\theta}$ do estado anterior e da entrada atual:

\begin{equation}
    h_t = f_{\theta}(h_{t-1}, x_t)
\end{equation}

Onde:

\begin{itemize}
    \item $h_t$ é o estado atual
    \item $f_{\theta}$ é uma função fixa com parâmetros $\theta$
    \item $h_{t-1}$ é o estado anterior
    \item $x_t$ é o vetor de entrada atual
\end{itemize}

Em mais detalhes, podemos aprofundar a rede como:


\begin{equation}
    h_t = tanh(W_{hh}h_{t-1}+W_{xh}x_{t}+b_h)
\end{equation}
Onde:
\begin{itemize}
    \item $W_{hh}$ É a matriz de pesos que conecta o estado oculto anterior ($h_{t-1}$) ao estado oculto atual ($h_t$). Esses pesos ditam o quão fortemente a informação do passado influencia o presente.
    \item $W_{xh}$ É a matriz de pesos que conecta a entrada atual ($x_t$) ao estado oculto atual ($h_t$). Esses pesos determinam a importância da entrada corrente na atualização do estado.
    \item $b_h$ é o viés para estado oculto. Ele permite que a rede aprenda um deslocamento na ativação.
    \item $tanh(^.)$ é a função de ativação não-linear.
\end{itemize}


Além do estado oculto, uma RNN também pode produzir uma saída ($y_t$) em cada passo de tempo.  Essa saída também é tipicamente uma função do estado oculto atual:

\begin{equation}
    y_t = W_{hy}h_t + b_y
\end{equation}

Onde:
\begin{itemize}
    \item $W_{hy}$ É a matriz de pesos que mapeia o estado oculto atual ($h_t$) para a saída ($y_t$)
    \item $b_y$ É o vetor de bias para a saída.
\end{itemize}
\end{document}